%==============
\begin{center}
{\LARGE \bf Constraints Programming
}
\end{center}
%==============
\bigskip

{\bf Organizer}

\medskip

Akira Aiba (aiba@icot.or.jp)

\bigskip

{\bf Description}
\medskip

In these years, roles of constraints in problem solving getting more important and promising. On
the other hand, capabilities of computer algebra will increase problem-solving ability of
constraints. 

One of the combination of constraints and computer algebra is constraint logic programming
(CLP) with a constraint-solver having computer algebra related algorithm for constraint solving.
Groebner-base algorithm for CAL, GDCC and the quantifier elimination algorithm for
RISC-CLP are examples. 

In addition fot constraint solvers in CLP, there would be various interesting relations between
constraint solving/constraint handling and computer algebra. 

The main purpose of this session is to study and discuss the use of computer algebra and its related
technoques in constraint solving/constraint handling not restricted to CLP, but constraint
satisfaction including distributed constraint, propagation, relaxation, and so forth. 

\bigskip
{\bf Talks}
\medskip

Date: July 19th (Friday) 

\begin{description}
\item[8:30-09:00] \ \\
  Yosuke Sato\\
  {\bf Application of Groebner basis in constraint of non-numerical domains.} 
\item[09:00-09:30] \ \\
  Eric Monfroy\\
  {\bf Solver collaborations for non-linear constraints.
} 
\item[09:30-10:00] \ \\
  Jason Harris\\
  {\bf Semantic Pattern Matching In Mathematica.
} 
\item[10:00-10:30] \ \\
  \ \\
  {\bf Coffee Break}
\item[10:30-11:00] \ \\
  Olga Caprotti, Hoon Hong\\
  {\bf Solving Equational Constraints via Patterns.} 
\item[11:00-11:30] \ \\
  Frederic Benhamou, Laurent Granvilliers
\\
  {\bf Application of Groebner Bases for Numerical Constraint Solving.
} 
\item[11:30-12:00] \ \\
  Hoon Hong, Stefan Ratschan\\
  {\bf Determining the Relationship Among Sets.
} 
\end{description}

